% arXiv submission: candle-cli
% Compile: pdflatex paper.tex
\documentclass[11pt]{article}

\usepackage[utf8]{inputenc}
\usepackage[T1]{fontenc}
\usepackage[margin=1in]{geometry}
\usepackage[hidelinks]{hyperref}
\usepackage{graphicx}
\usepackage{booktabs}
\usepackage{listings}
\usepackage{xcolor}
\usepackage{enumitem}
\usepackage{amsmath}
\usepackage{amssymb}

\definecolor{codebg}{rgb}{0.95,0.95,0.95}
\lstset{
    basicstyle=\ttfamily\small,
    backgroundcolor=\color{codebg},
    frame=single,
    breaklines=true,
    numbers=none,
    aboveskip=6pt,
    belowskip=6pt,
}

\title{\textbf{Candle-CLI: A Lightweight Agentic Coding Assistant}\\
\large with Multi-Agent Coordination, Layered Memory, and Sandboxed Execution}

% ============================================================
% FILL IN YOUR NAME AND EMAIL BEFORE SUBMITTING TO arXiv
% ============================================================
\author{Your Name\\
\texttt{your.email@example.com}}

\date{}

\begin{document}
\maketitle

\begin{abstract}
We present Candle-CLI, an open-source agentic coding assistant built with a Rust core and Python bridge architecture. Inspired by the lightweight AI runtime philosophy of the Candle project, Candle-CLI implements a bounded multi-step agent loop with text-based JSON tool calling, enabling large language models to read, search, edit, and execute code within local development environments. The system features eight built-in tools, a four-mode permission system with workspace path boundary enforcement, Docker-based sandboxed shell execution, multi-agent task delegation via sub-agent dispatch, and a layered memory architecture combining session persistence with project-level persistent knowledge. We introduce grep-based retrieval-augmented pre-search to reduce API round trips, fault tolerance with exponential backoff retry, and an observability stack comprising tool listing, runtime status inspection, and execution trace with millisecond-level step timing and structured JSON export. A scenario harness framework enables automated end-to-end evaluation across multiple model backends. Candle-CLI adopts an API-first strategy supporting DeepSeek, Ollama, vLLM, and OpenAI through a persistent bridge worker with streaming token output. The system is implemented in approximately 4,800 lines of Rust and Python, with 85 passing tests, and is freely available on GitHub at \url{https://github.com/DuangZ-GR/candle-cli}.
\end{abstract}

%---------------------------------------------------------------------
\section{Introduction}

Large language models (LLMs) have demonstrated remarkable capabilities in code generation and software engineering tasks. However, most LLM-based coding tools operate within constrained environments---either as browser-based chat interfaces or IDE plugins with limited access to the underlying development workflow. These systems typically lack the ability to autonomously read project files, search codebases, execute shell commands, and edit source code in a controlled, observable manner.

Agentic coding assistants address this gap by equipping LLMs with tool-calling capabilities, enabling them to interact with local development environments through structured action requests. Existing systems such as Claude Code, Aider, SWE-Agent, and OpenHands have demonstrated the potential of this approach, but they often rely on proprietary APIs, complex dependency stacks, or design choices that limit extensibility and local deployment.

We introduce Candle-CLI, a lightweight, open-source agentic coding assistant designed with the following principles:

\begin{enumerate}[leftmargin=*]
    \item \textbf{Rust-first reliability.} The control plane---CLI interaction, agent loop orchestration, tool execution, permission enforcement, session management, and observability---is implemented in Rust, providing compile-time type safety and predictable runtime behavior.
    \item \textbf{API-first accessibility.} Model inference uses OpenAI-compatible APIs by default, eliminating GPU requirements. A persistent Python bridge worker handles backend adaptation, supporting DeepSeek, Ollama, vLLM, and OpenAI with zero-code switching via environment variables.
    \item \textbf{Multi-agent coordination.} A task delegation tool enables the main agent to spawn read-only sub-agents with bounded execution loops for isolated subtask handling.
    \item \textbf{Layered memory.} Beyond session-level conversation persistence, the system maintains project-level memory stored in structured JSON, automatically injected into the system prompt across sessions.
    \item \textbf{Sandboxed execution.} Shell commands can be optionally executed inside Docker containers with read-only workspace mounts and network isolation.
    \item \textbf{Built-in observability.} Three inspection interfaces---tool listing, runtime status, and execution tracing with millisecond timing and JSON export---provide comprehensive visibility into system behavior.
\end{enumerate}

The remainder of this paper is organized as follows. Section~\ref{sec:related} reviews related work. Section~\ref{sec:design} describes the system architecture and agent loop design. Section~\ref{sec:tools} details the tool system and permission architecture. Section~\ref{sec:advanced} covers advanced features including multi-agent coordination, layered memory, fault tolerance, and observability. Section~\ref{sec:evaluation} presents evaluation results. Section~\ref{sec:discussion} discusses limitations and future work, and Section~\ref{sec:conclusion} concludes.

%---------------------------------------------------------------------
\section{Related Work}
\label{sec:related}

\subsection{Agentic Coding Assistants}

\textbf{Claude Code}~\cite{claudecode} is Anthropic's official CLI coding agent, providing file reading, editing, shell execution, and multi-turn conversation with persistent sessions. It uses Anthropic's native tool-use API and supports sub-agent dispatch for parallel task execution. Candle-CLI shares a similar architectural philosophy but uses a text-based JSON tool-calling protocol that is model-agnostic rather than tied to a specific provider's function-calling API.

\textbf{Aider}~\cite{aider} is an open-source AI pair programming tool focused on editing workflows, supporting multiple model backends with precise code modifications and diff visualization. Candle-CLI complements this by focusing on the broader agentic infrastructure---permission control, multi-agent coordination, layered memory, and observability.

\textbf{SWE-Agent}~\cite{sweagent} and \textbf{OpenHands}~\cite{openhands} are research-oriented agentic coding frameworks targeting autonomous software engineering tasks such as GitHub issue resolution. They employ sophisticated agent-computer interfaces and multi-step reasoning. Candle-CLI is comparatively lightweight, designed as a terminal-native tool with an emphasis on deployability and extensibility.

\subsection{The Candle Project}

The Candle project~\cite{candle} is a pure-Python deep learning framework emphasizing lightweight deployment (approximately 10~MB install size), multi-hardware backend support (CPU, CUDA, Apple MPS, Ascend NPU), and PyTorch-compatible API design. Its roadmap outlines a four-phase progression from Foundation to Self-Hosting, with Cognitive Runtime and Agentic Kernel phases envisioning local model deployment integrated with self-debugging and self-improving agent systems. Candle-CLI draws architectural inspiration from this vision.

\subsection{Tool Calling Protocols}

LLM tool calling has evolved from proprietary API-specific formats (OpenAI function calling, Anthropic tool use) to more general approaches. Candle-CLI adopts a text-based JSON protocol using \texttt{<tool\_call>} tags with embedded JSON payloads. This approach is model-agnostic, compatible with any text-generation model capable of following formatting instructions, and does not require API-level tool-calling support.

%---------------------------------------------------------------------
\section{System Design}
\label{sec:design}

\subsection{Architecture Overview}

Candle-CLI employs a layered architecture with clear separation between the system control plane (Rust) and the model inference layer (Python bridge).

\begin{center}
\begin{minipage}{0.45\textwidth}
\footnotesize\ttfamily
\begin{tabbing}
xxx\=xxxxxxxxxxxx\=\kill
CLI Layer          \> (src/cli/)   \> args, repl, cmds\\
Agent Layer        \> (src/agent/) \> loop, tool\_call, trace\\
                   \>              \> $\to$ Permission Layer\\
Model Layer        \> (src/model/) \> runtime, mock, bridge, candle\\
Tool Layer         \> (src/tools/) \> registry, builtin (8 tools)\\
Session Layer      \> (src/session/)\> model, store, resume, memory\\
                   \>              \> $\to$ Permission (mode, policy)\\
Python Bridge      \> (python/)    \> worker, runtime, prompt, config\\
LLM Backends       \>              \> DeepSeek, Ollama, vLLM, OpenAI
\end{tabbing}
\end{minipage}
\end{center}

The \textbf{CLI Layer} handles command-line argument parsing, REPL interaction loop (with readline editing), and slash command dispatch. The \textbf{Agent Layer} implements the bounded multi-step agent loop, text JSON tool-call parser, execution trace recorder, and turn finalization. The \textbf{Model Layer} defines the \texttt{CandleTargetRuntime} trait with three implementations: \texttt{MockRuntime} (returns stub responses for testing), \texttt{LocalBridgeRuntime} (communicates with a persistent Python subprocess via stdin/stdout JSON line protocol), and \texttt{CandleRuntime} (placeholder for future local inference). The \textbf{Tool Layer} manages tool registration, parameter parsing, workspace path boundary enforcement, and execution dispatch. The \textbf{Permission Layer} defines four permission modes and enforces access control at the agent loop level. The \textbf{Session Layer} handles conversation context serialization to JSON files and project-level persistent memory.

\subsection{Agent Loop Design}

The agent loop is the core execution engine, shown in Algorithm~1.

\begin{lstlisting}[caption={Agent loop (simplified Rust pseudocode)}, label={alg:loop}]
fn run_agent_loop(session, runtime, tools, policy, max_steps=8):
    for step in 0..max_steps:
        request = build_turn_request(session, tools_json)
        result = runtime.generate_turn(request)

        match parse_tool_call(result.text):
            Some(tool_call):
                if !policy.allows(tool_call.name):
                    record_error("not allowed in current mode"); continue
                if policy.requires_prompt(tool_call.name)
                   && !confirm_dangerous_action(...):
                    record_error("denied by user"); continue
                output = tools.execute(tool_call.name, tool_call.input)
                record_result(output); continue
            None:
                return result.text          // final answer
            Err(parse_error):
                append_correction(parse_error); continue
    return "max steps reached"
\end{lstlisting}

The loop has three key properties: (1) it is bounded (default 8 steps), preventing infinite tool-calling; (2) parse errors are recoverable---malformed tool calls trigger corrective guidance rather than termination; and (3) permission denials are communicated back to the model as structured error messages, enabling strategy adjustment.

\subsection{Tool Call Protocol}

Candle-CLI uses a text-based JSON protocol. Models emit requests as:

\begin{lstlisting}[caption={Tool call format}, label={lst:tc}]
<tool_call>{"id":"call-1","name":"read",
 "input":{"file_path":"README.md"}}</tool_call>
\end{lstlisting}

The parser searches for \texttt{<tool\_call>} tags, extracts the enclosed JSON, and validates three required fields: \texttt{id} (string), \texttt{name} (string), and \texttt{input} (object). Six error types are defined, each with descriptive recovery guidance. A fallback parser also accepts function-style calls: \texttt{read(\{"file\_path": "README.md"\})}. This protocol is model-agnostic---compatible with any text-generation model capable of following formatting instructions.

\subsection{Python Bridge and API-First Strategy}

The Python bridge runs as a persistent subprocess, communicating with Rust via JSON line protocol over stdin/stdout. Key design decisions:

\begin{itemize}[leftmargin=*]
    \item \textbf{Persistent worker.} The bridge remains alive across turns, eliminating cold-start overhead.
    \item \textbf{Streaming output.} API requests use \texttt{"stream": true}, enabling real-time token display during generation.
    \item \textbf{Backend adaptation.} \texttt{bridge\_prompt.py} converts Rust session messages (including \texttt{ToolCall} and \texttt{ToolResult} content blocks) into the format expected by the target API.
    \item \textbf{API-first rationale.} Requiring only environment variables for configuration, this approach eliminates GPU requirements while preserving the option to integrate local inference backends in the future through the \texttt{CandleTargetRuntime} trait.
\end{itemize}

%---------------------------------------------------------------------
\section{Tool System and Permission Architecture}
\label{sec:tools}

\subsection{Built-in Tools}

Candle-CLI provides eight built-in tools, summarized in Table~\ref{tab:tools}.

\begin{table}[h]
\centering
\caption{Built-in tools in Candle-CLI.}
\label{tab:tools}
\begin{tabular}{@{}llll@{}}
\toprule
\textbf{Tool} & \textbf{Input} & \textbf{Purpose} & \textbf{Mutation} \\
\midrule
\texttt{pwd}       & \texttt{\{\}}                  & Show workspace directory        & No \\
\texttt{read}      & \texttt{\{"file\_path":"..."\}}  & Read UTF-8 file (path enforced) & No \\
\texttt{glob}      & \texttt{\{"pattern":"..."\}}     & Find files by glob pattern      & No \\
\texttt{grep}      & \texttt{\{"pattern":"...","path":"..."\}} & Search file contents     & No \\
\texttt{web\_search}& \texttt{\{"query":"..."\}}       & Web search (dual-backend)       & No \\
\texttt{task}      & \texttt{\{"description":"..."\}}  & Delegate to read-only sub-agent & No \\
\texttt{edit}      & \texttt{\{"file\_path":"...",...\}}& Exact-once text replacement     & Yes \\
\texttt{shell}     & \texttt{\{"command":"..."\}}      & Run shell command (w/ timeout)  & Possible \\
\bottomrule
\end{tabular}
\end{table}

\textbf{Read tools} (\texttt{pwd}, \texttt{read}, \texttt{glob}, \texttt{grep}) provide codebase inspection. \texttt{read} enforces workspace path boundary checks. \texttt{grep} performs recursive text search with \texttt{path:line:content} output. \texttt{glob} supports \texttt{*.rs} and \texttt{**/*.rs} patterns.

\textbf{Network tool} (\texttt{web\_search}) uses DuckDuckGo Lite with automatic fallback to Sogou. HTML cleaning employs Python \texttt{re.DOTALL} mode for multi-line script/style block removal.

\textbf{Coordination tool} (\texttt{task}) spawns read-only sub-agents with a 3-step bounded loop for isolated subtask handling (Section~\ref{sec:multiagent}).

\textbf{Mutation tools} (\texttt{edit}, \texttt{shell}) perform file modifications and command execution. \texttt{edit} uses exact-once match semantics, failing on zero or multiple matches. \texttt{shell} supports configurable timeout with SIGKILL and structured output capturing exit codes.

\subsection{Permission System}

The permission system provides four execution modes configured via \texttt{CANDLE\_CLI\_PERMISSION} (Table~\ref{tab:perms}).

\begin{table}[h]
\centering
\caption{Permission modes.}
\label{tab:perms}
\begin{tabular}{@{}ll@{}}
\toprule
\textbf{Mode} & \textbf{Behavior} \\
\midrule
\texttt{read-only}          & Only \texttt{pwd}, \texttt{read}, \texttt{glob}, \texttt{grep} allowed \\
\texttt{workspace-write}    & All tools allowed without confirmation (default) \\
\texttt{prompt}             & Read tools auto-allowed; mutation tools require interactive confirmation \\
\texttt{danger-full-access} & All tools allowed without confirmation \\
\bottomrule
\end{tabular}
\end{table}

Permission checks are embedded in the agent loop's execution path, not in a separate interceptor layer. Workspace boundary enforcement (\texttt{resolve\_existing\_path}, \texttt{resolve\_writable\_path}) is implemented at the \texttt{ToolRegistry} level. On Unix systems, symlink escape detection prevents operations from following symbolic links pointing outside the workspace. The \texttt{prompt} mode supports both interactive terminal confirmation and programmatic responses via the \texttt{CANDLE\_CLI\_PERMISSION\_RESPONSE} environment variable.

%---------------------------------------------------------------------
\section{Advanced Features}
\label{sec:advanced}

\subsection{Multi-Agent Coordination}
\label{sec:multiagent}

The \texttt{task} tool implements a delegation-based multi-agent coordination pattern. When invoked: (1) a fresh sub-session is created with only the task description; (2) a read-only tool registry is instantiated; (3) a bounded agent loop runs with a maximum of 3 steps; and (4) the final answer is returned to the main agent as a structured tool result. Sub-agents operate with minimal privileges and bounded scope, preventing cascading failures or unauthorized modifications.

\subsection{Layered Memory Architecture}

Candle-CLI implements a two-tier memory system:

\textbf{Session memory} (short-term): Conversation history is serialized as JSON and persisted after each turn. Sessions can be listed, resumed, and cleared.

\textbf{Project memory} (long-term): A structured JSON file (\texttt{.candle-cli/memory.json}) stores three categories: (1) key file paths, (2) commonly used commands, and (3) free-form notes. Memory is loaded at startup and injected into the system prompt on every turn. Users manage memory via \texttt{/memory file}, \texttt{/memory cmd}, and \texttt{/memory note} subcommands.

\subsection{Grep-Based RAG Pre-Search}

Before each model call, the context builder performs a lightweight retrieval step: extracts the last user message, filters out non-technical chat, derives search phrases (full question for short queries, component phrases split by delimiters for longer ones), executes \texttt{grep} against \texttt{src/}, and injects up to 10 matching lines into the user message as context. This reduces API round trips by providing relevant code context before the model's first reasoning step.

\subsection{Fault Tolerance}

\textbf{API retry with exponential backoff.} The Python bridge retries failed API calls up to three times with delays of 1s, 2s, and 4s. HTTP 4xx errors are not retried.

\textbf{Shell timeout.} Commands timeout after a configurable duration (default 30 seconds) with SIGKILL enforcement.

\textbf{Tool call error recovery.} Malformed tool calls produce descriptive corrective messages. The loop continues, giving the model an opportunity to retry.

\subsection{Observability Stack}

Three inspection commands provide transparency:

\texttt{/tools}: Lists all registered tools, showing the system's current capability boundary.

\texttt{/status}: Displays runtime snapshot including session ID, message count with token estimation, workspace path, permission mode, runtime type, model identifier, and configured turn limit.

\texttt{/trace}: Shows the last execution chain with per-step millisecond timing and JSON export. The \texttt{ExecutionTrace} records six event types: \texttt{BuildTurnRequest}, \texttt{RuntimeGenerateTurn}, \texttt{ParseToolCall}, \texttt{ToolCall}, \texttt{ToolResult}, and \texttt{FinalAnswer}. It captures system-level execution events, not model-internal reasoning.

\subsection{Docker Sandbox}

Setting \texttt{CANDLE\_CLI\_SANDBOX=docker} executes shell commands inside an Alpine Linux container with the workspace mounted read-only and network access disabled (\texttt{--network=none}). This provides defense-in-depth for scenarios where model-generated commands should not have filesystem write access or network connectivity.

%---------------------------------------------------------------------
\section{Evaluation}
\label{sec:evaluation}

\subsection{Harness Framework}

Candle-CLI includes a built-in scenario harness (\texttt{cargo run -- harness}) that evaluates four predefined tasks: reading a project file, performing a glob search, searching source code with grep, and executing a shell command. Each scenario is evaluated on pass/fail status, elapsed time, tool execution steps, API call rounds, and estimated token consumption. Results are exported as structured JSON.

\subsection{Codebase}

The system comprises approximately 4,800 lines of code across 64 source files, with 85 Rust tests and 31 Python tests passing. The Rust core (\texttt{src/}, 40 files) implements all system control logic. The Python bridge (\texttt{python/}, 6 files) handles model backend adaptation.

%---------------------------------------------------------------------
\section{Discussion and Limitations}
\label{sec:discussion}

\textbf{Model dependency.} Effectiveness depends on the underlying model's ability to follow the tool-calling protocol. Models below approximately 3B parameters may produce malformed tool calls. The system mitigates this through format error recovery, but fundamental quality depends on model capability.

\textbf{API dependency.} The API-first strategy simplifies deployment but introduces latency and cost. The persistent bridge worker reduces per-turn overhead. Future integration with the Candle framework for local inference would address this.

\textbf{Security scope.} Workspace path enforcement and the Docker sandbox provide meaningful protection but are not a complete security solution. The Docker sandbox requires Docker installation and explicit opt-in. The permission system provides defense-in-depth and should be supplemented with OS-level isolation for production deployments.

\textbf{Evaluation scope.} The current harness covers four basic scenarios. Comprehensive evaluation on established benchmarks (such as SWE-bench~\cite{sweagent} or HumanEval) would provide more rigorous assessment.

%---------------------------------------------------------------------
\section{Conclusion and Future Work}
\label{sec:conclusion}

We have presented Candle-CLI, an open-source agentic coding assistant combining a Rust-based control plane with a Python bridge for model backend adaptation. The system implements a bounded multi-step agent loop, eight development tools, a four-mode permission architecture, multi-agent task delegation, layered memory, grep-based RAG pre-search, and an observability stack with execution tracing and millisecond timing.

Future work includes: (1) integrating the Candle framework as a local inference backend for hybrid deployment; (2) implementing patch-style editing with diff preview; (3) expanding evaluation to established coding benchmarks; (4) adding MCP (Model Context Protocol) support for external tool integration; and (5) developing a self-development workflow where Candle-CLI modifies, tests, and improves its own codebase.

The source code is available at \url{https://github.com/DuangZ-GR/candle-cli}.

%---------------------------------------------------------------------
\section*{Acknowledgments}

The authors thank the Candle project for its inspiring vision of lightweight, agentic-ready AI runtime infrastructure, and the open-source communities behind Rust, Python, DeepSeek, Ollama, and the broader LLM ecosystem.

\bibliographystyle{plain}
\begin{thebibliography}{10}

\bibitem{candle}
Candle Project. \textit{Candle: A pure-Python deep learning framework.} \url{https://github.com/candle-org/candle}, 2024.

\bibitem{claudecode}
Anthropic. \textit{Claude Code: Agentic coding tool for the terminal.} \url{https://docs.anthropic.com/en/docs/claude-code}, 2025.

\bibitem{aider}
P.~Gauthier. \textit{Aider: AI pair programming in your terminal.} \url{https://github.com/Aider-AI/aider}, 2024.

\bibitem{sweagent}
J.~Yang, C.~Jimenez, A.~Wettig, et al. \textit{SWE-Agent: Agent-Computer Interfaces Enable Automated Software Engineering.} arXiv:2405.15793, 2024.

\bibitem{openhands}
All Hands AI. \textit{OpenHands: Code Less, Make More.} \url{https://github.com/All-Hands-AI/OpenHands}, 2024.

\bibitem{deepseek}
DeepSeek AI. \textit{DeepSeek-V4: Advancing Open-Source Language Models.} \url{https://api.deepseek.com}, 2025.

\bibitem{rust}
S.~Klabnik, C.~Nichols. \textit{The Rust Programming Language, 2nd Edition.} No Starch Press, 2023.

\bibitem{transformers}
T.~Wolf, L.~Debut, V.~Sanh, et al. \textit{HuggingFace Transformers: State-of-the-art Natural Language Processing.} In \textit{EMNLP}, 2020.

\end{thebibliography}

\end{document}
